problem:  
Let \((X,L)\) be a polarized compact Kähler manifold, and for \(1<p<\infty\), let \((\mathcal E^p,d_p)\) denote the finite‑energy completion of the space of Kähler potentials endowed with the standard \(L^p\) Finsler metric  
\[
\|\dot\varphi\|_{p,\varphi}=\Big(\int_X |\dot\varphi|^p\,\omega_\varphi^n\Big)^{1/p}.
\]
For each \(p\), consider the Mabuchi functional \(E:\mathcal E^p\to(-\infty,+\infty]\), and let its **non‑Archimedean counterpart** \(E^{\mathrm{NA}}\) be defined on the space of test configurations of \((X,L)\).  It is known that along the geodesic ray associated to a test configuration \((\mathcal X,\mathcal L)\),
\[
E(\varphi_t)=tE^{\mathrm{NA}}(\mathcal X,\mathcal L)+o(t),
\]
with the weak geodesics governed by the complex Monge–Ampère equation.  

Define for each \(p>1\) the *metric slope function* of \(E\) at a point \(\varphi\in\mathcal E^p\):
\[
|\nabla E|_{p}(\varphi):=\limsup_{\psi\to\varphi}\frac{[E(\varphi)-E(\psi)]_+}{d_{p}(\varphi,\psi)}.
\]
  
**Question.**
1. For uniformly K‑stable \((X,L)\), is the slope function \(|\nabla E|_{p}(\varphi)\) bounded below (in an appropriate weak sense) only in the \(p\)-range where the Mabuchi functional is \(p\)-coercive? Could such quantitative lower bounds interpolate between uniform K‑stability (\(p=2\)) and non‑Archimedean coercivity as \(p\to\infty\)?  
2. Does the large‑\(p\) asymptotic behavior of \(|\nabla E|_p\) along rays associated to test configurations detect new stability thresholds or new energy invariants beyond \(E^{\mathrm{NA}}\)?  
3. If one defines the **\(p\)-gradient flow** of \(E\) by the metric‑slope flow on \((\mathcal E^p,d_p)\), does the limiting flow as \(p\to1\) or \(p\to\infty\) produce novel pluripotential or non‑Archimedean dynamics reflecting algebraic degeneration?  

Why is it a "good" problem:  
This formulation centers on a rigorous, well‑defined geometric analytic object—the metric slope of the Mabuchi functional on \((\mathcal E^p,d_p)\)—whose behavior as \(p\) varies has never been analyzed systematically.  Unlike the always‑known convexity of \(E\), the slope quantitatively measures coercivity and flow dynamics in the \(L^p\) metric geometry, offering a unified analytic link between the finite‑energy pluripotential theory and non‑Archimedean energy functionals.  Studying its \(p\)-dependence and asymptotics could uncover a new family of stability invariants that interpolate between differential‑geometric and algebraic notions of stability and provide a metric interpretation of uniform K‑stability.  The statement is precise, logically sound, and conceptually simple, yet its resolution would require new interaction between gradient‑flow theory in infinite‑dimensional Finsler spaces, pluripotential analysis, and the algebraic geometry of test configurations—meeting the standards of a profound and forward‑looking research problem.